%==================================================================
% PeriX -- Reference Manual
% Peridynamic Differential Operator framework for multiphysics
% Author: Yang Bai / MMLab
% License: GNU GPLv3
%==================================================================
\documentclass[11pt,a4paper]{article}

\usepackage[T1]{fontenc}
\usepackage{lmodern}
\usepackage[utf8]{inputenc}
\usepackage[english]{babel}
\usepackage{amsmath,amssymb,amsfonts}
\usepackage{mathtools}
\usepackage{geometry}
\geometry{a4paper,margin=2.4cm}
\emergencystretch=3em % a little extra inter-word stretch to avoid overfull lines
\usepackage{xcolor}
\usepackage{hyperref}
\hypersetup{
  colorlinks=true,
  linkcolor=blue!50!black,
  urlcolor=blue!60!black,
  citecolor=blue!50!black,
  bookmarksnumbered=true,
  pdftitle={PeriX Reference Manual},
  pdfauthor={Yang Bai}
}
\usepackage{listings}
\usepackage{booktabs}
\usepackage{tabularx}
\usepackage{enumitem}
\usepackage{microtype}
\usepackage{fancyhdr}
\usepackage{titling}
\usepackage{caption}
\usepackage{graphicx}

% --- Code listings (JSON / Bash / C++) ---
\definecolor{lstbg}{rgb}{0.97,0.97,0.97}
\definecolor{lstkw}{rgb}{0.10,0.30,0.70}
\definecolor{lstst}{rgb}{0.60,0.00,0.00}
\definecolor{lstcm}{rgb}{0.30,0.50,0.30}
\lstdefinelanguage{json}{
  morestring=[b]",
  morecomment=[l]{//},
  literate=
   *{0}{{{\color{lstst}0}}}{1}
    {1}{{{\color{lstst}1}}}{1}
    {2}{{{\color{lstst}2}}}{1}
    {3}{{{\color{lstst}3}}}{1}
    {4}{{{\color{lstst}4}}}{1}
    {5}{{{\color{lstst}5}}}{1}
    {6}{{{\color{lstst}6}}}{1}
    {7}{{{\color{lstst}7}}}{1}
    {8}{{{\color{lstst}8}}}{1}
    {9}{{{\color{lstst}9}}}{1}
    {:}{{{\color{lstkw}{:}}}}{1}
    {,}{{{\color{lstkw}{,}}}}{1}
    {\{}{{{\color{lstkw}{\{}}}}{1}
    {\}}{{{\color{lstkw}{\}}}}}{1}
    {[}{{{\color{lstkw}{[}}}}{1}
    {]}{{{\color{lstkw}{]}}}}{1}
}
\lstset{
  basicstyle=\ttfamily\footnotesize,
  backgroundcolor=\color{lstbg},
  keywordstyle=\color{lstkw}\bfseries,
  stringstyle=\color{lstst},
  commentstyle=\color{lstcm}\itshape,
  breaklines=true,
  breakatwhitespace=true,
  columns=fullflexible,
  showstringspaces=false,
  frame=single,
  framerule=0.3pt,
  rulecolor=\color{black!30},
  tabsize=2,
  captionpos=b,
  xleftmargin=8pt,
  xrightmargin=4pt,
  literate=
    {Pe}{Pe}{2}
}

% Inline code that may break after the usual separators (/ . : -) so long
% file paths and qualified names wrap instead of overrunning the right margin.
% The argument is re-scanned with those characters made active; \_ , spaces and
% other LaTeX markup are preserved. Robust so it also works in section titles.
\makeatletter
% Break opportunities for the re-scan below: preferred (zero-penalty) breaks
% after path / qualifier separators "/ . : -", and fall-back breaks before an
% upper-case letter (camelCase) that are taken only when no separator break
% keeps the line inside the margin. "=" is intentionally NOT made active (it
% would corrupt the \catcode...=\active assignments themselves).
\begingroup
\catcode`\/=\active \catcode`\.=\active \catcode`\:=\active \catcode`\-=\active
\catcode`\A=\active \catcode`\B=\active \catcode`\C=\active \catcode`\D=\active
\catcode`\E=\active \catcode`\F=\active \catcode`\G=\active \catcode`\H=\active
\catcode`\I=\active \catcode`\J=\active \catcode`\K=\active \catcode`\L=\active
\catcode`\M=\active \catcode`\N=\active \catcode`\O=\active \catcode`\P=\active
\catcode`\Q=\active \catcode`\R=\active \catcode`\S=\active \catcode`\T=\active
\catcode`\U=\active \catcode`\V=\active \catcode`\W=\active \catcode`\X=\active
\catcode`\Y=\active \catcode`\Z=\active
\gdef\perix@codebrk{%
  \def/{\char`\/\allowbreak}\def.{\char`\.\allowbreak}%
  \def:{\char`\:\allowbreak}\def-{\char`\-\allowbreak}%
  \def\_{\textunderscore\allowbreak}% break after underscores too (file names)
  \def A{\penalty200\char`\A}\def B{\penalty200\char`\B}\def C{\penalty200\char`\C}%
  \def D{\penalty200\char`\D}\def E{\penalty200\char`\E}\def F{\penalty200\char`\F}%
  \def G{\penalty200\char`\G}\def H{\penalty200\char`\H}\def I{\penalty200\char`\I}%
  \def J{\penalty200\char`\J}\def K{\penalty200\char`\K}\def L{\penalty200\char`\L}%
  \def M{\penalty200\char`\M}\def N{\penalty200\char`\N}\def O{\penalty200\char`\O}%
  \def P{\penalty200\char`\P}\def Q{\penalty200\char`\Q}\def R{\penalty200\char`\R}%
  \def S{\penalty200\char`\S}\def T{\penalty200\char`\T}\def U{\penalty200\char`\U}%
  \def V{\penalty200\char`\V}\def W{\penalty200\char`\W}\def X{\penalty200\char`\X}%
  \def Y{\penalty200\char`\Y}\def Z{\penalty200\char`\Z}%
}
\endgroup
\DeclareRobustCommand{\code}[1]{%
  \ifmmode
    \text{\ttfamily\small #1}% in math mode: plain monospace, no re-scan
  \else
    \begingroup
      \ttfamily\small
      \catcode`\/=\active \catcode`\.=\active \catcode`\:=\active \catcode`\-=\active
      \catcode`\A=\active \catcode`\B=\active \catcode`\C=\active \catcode`\D=\active
      \catcode`\E=\active \catcode`\F=\active \catcode`\G=\active \catcode`\H=\active
      \catcode`\I=\active \catcode`\J=\active \catcode`\K=\active \catcode`\L=\active
      \catcode`\M=\active \catcode`\N=\active \catcode`\O=\active \catcode`\P=\active
      \catcode`\Q=\active \catcode`\R=\active \catcode`\S=\active \catcode`\T=\active
      \catcode`\U=\active \catcode`\V=\active \catcode`\W=\active \catcode`\X=\active
      \catcode`\Y=\active \catcode`\Z=\active
      \perix@codebrk
      \endlinechar=\m@ne
      \scantokens{#1}%
    \endgroup
  \fi}
\makeatother
% In PDF bookmarks / metadata the breakable form cannot expand; use plain text.
\pdfstringdefDisableCommands{\def\code#1{#1}}
\newcommand{\PeriX}{\textsf{PeriX}}
\newcommand{\Lap}{\nabla^{2}}
\newcommand{\Grad}{\nabla}
\newcommand{\Div}{\nabla\cdot}

\title{
  \vspace{-1.4cm}
  \textbf{\textsf{PeriX} -- Reference Manual} \\[0.4em]
  \large Peridynamic Differential Operator (PDDO) Framework for Multiphysics Simulation
}
\author{Yang Bai\\\small Materials Mechanics Lab (MMLab) \\ \small \texttt{yangbai90@outlook.com}}
\date{Release tag: 2026-04-16, version 0.10}

\pagestyle{fancy}
\fancyhf{}
\fancyhead[L]{\PeriX{} Manual}
\fancyhead[R]{\thepage}

%==================================================================
\begin{document}
\maketitle
\thispagestyle{empty}

\begin{abstract}
\noindent
\PeriX{} is an open-source C++20 framework for solving steady and
transient multiphysics problems by means of the
\emph{Peridynamic Differential Operator} (PDDO). The PDDO is a
nonlocal collocation-style method due to E.~Madenci and co-workers
that reproduces ordinary partial derivatives of arbitrary order on
a point cloud inside a finite \emph{horizon} $\delta$, using a
weighted least-squares fit against a polynomial basis. \PeriX{}
ships with seven physics kernels -- Poisson, transient diffusion,
Cahn--Hilliard phase separation, implicit and explicit bond-breaking
peridynamic fracture, and stress-assisted Cahn--Hilliard transport
with and without fracture -- driven entirely by a single JSON input
file. This manual documents the public release: the build system,
the command line, the complete JSON input schema with every
recognised key, its type, default value and validation rule, the
mathematical formulation of each kernel exactly as implemented, the
boundary- and initial-condition machinery, the linear and nonlinear
solvers, the output system, and all bundled example problems.
Every statement in this manual has been cross-checked against the
source code of the release.
\end{abstract}

\vspace{0.5em}
\begin{center}
\begin{tabular}{ll}
\toprule
Repository & \url{https://github.com/MatMechLab/PeriX} \\
License    & GNU General Public License v3.0 (GPLv3) \\
Language   & C++20 (CMake $\geq$ 3.21) \\
Input      & a single JSON file (\code{perix -i input.json}) \\
Output     & VTK \code{.vtu}/\code{.pvd} and/or ExodusII \code{.e} \\
\bottomrule
\end{tabular}
\end{center}

\newpage
\tableofcontents
\newpage

%==================================================================
\section{Introduction}
%==================================================================

\subsection{What is \PeriX{}?}

\PeriX{} discretises partial differential equations on a
\emph{peridynamic point cloud}: every material point $i$ carries a
position $\mathbf{x}_i$, a volume $V_i$, and a family
$\mathcal{H}_i$ of neighbouring points within a horizon $\delta$.
Spatial derivatives are replaced by the Peridynamic Differential
Operator, a set of per-bond weights $g_i^{\,q}(\mathbf{x}_j)$ such
that, for a smooth field $u$,
\begin{equation}
  \left.\frac{\partial^{|q|} u}{\partial x^{q_1}\partial y^{q_2}\partial z^{q_3}}\right|_{\mathbf{x}_i}
  \;\approx\; \sum_{j\in\mathcal{H}_i} g_i^{\,q}(\mathbf{x}_j)\,
  \bigl(u_j-u_i\bigr),
  \label{eq:pddo-idea}
\end{equation}
where the neighbour volume $V_j$ is already contained in
$g_i^{\,q}$. Because \eqref{eq:pddo-idea} needs no mesh topology,
no shape functions and no integration rule, the same machinery
serves diffusion-type equations, fourth-order Cahn--Hilliard
systems, and bond-based dynamic fracture on one code path.

The public release contains seven ready-to-use physics kernels
(Table~\ref{tab:kernels}); each is selected by a single
\code{type} string in the \code{ElmtSystem} block of the input
file.

\begin{table}[htbp]
\centering
\caption{Physics kernels in the public release.}
\label{tab:kernels}
\footnotesize
\begin{tabularx}{\textwidth}{llX}
\toprule
\textbf{Kernel \code{type}} & \textbf{DoFs} & \textbf{Physics} \\
\midrule
\code{poisson}          & \code{u}  & steady Poisson equation $\sigma\Lap u + f = 0$ \\
\code{diffusion}        & \code{c}  & transient diffusion $\dot c = D\Lap c + f$ (backward Euler) \\
\code{cahnhilliard}     & \code{c},\code{mu} & Cahn--Hilliard phase separation, mixed two-field split \\
\code{pddo\_dynamic\_frac} & \code{ux},\code{uy}(,\code{uz}) & implicit PDDO elastodynamics with critical-stretch bond failure \\
\code{explicit\_pddo\_frac} & \code{ux},\code{uy}(,\code{uz}) & explicit (central-difference) matrix-free variant of the above \\
\code{stress\_cahnhilliard} & \code{c},\code{mu},\code{ux},\code{uy}(,\code{uz}) & stress-assisted Cahn--Hilliard transport coupled to elasticity \\
\code{frac\_stress\_cahnhilliard} & \code{c},\code{mu},\code{ux},\code{uy}(,\code{uz}) & stress-assisted CH with bond-breaking fracture \\
\bottomrule
\end{tabularx}
\end{table}

\subsection{Scope of this manual}

This manual documents exactly what the public repository accepts
and computes. The JSON validator
(\code{src/InputSystem/InputSchema.cpp}) is strict: any unknown
block, unknown key, wrong type or out-of-range value aborts the run
with a descriptive error message before any computation starts.
Consequently, every key listed in Section~\ref{sec:input} is the
complete public vocabulary -- there are no hidden options.

\subsection{Source layout}

\begin{table}[htbp]
\centering
\caption{Repository layout.}
\label{tab:layout}
\footnotesize
\begin{tabularx}{\textwidth}{lX}
\toprule
\textbf{Directory} & \textbf{Contents} \\
\midrule
\code{src/}, \code{include/} & C++ sources and headers, one directory per subsystem (\code{Mesh}, \code{PDMesh}, \code{ElmtSystem}, \code{BCSystem}, \code{ICSystem}, \code{EquationSystem}, \code{LinearSolver}, \code{NonlinearSolver}, \code{TimeStepping}, \code{OutputSystem}, \code{InputSystem}, \code{PDProblem}, \code{Utils}) \\
\code{external/} & vendored header-only dependencies: Eigen and nlohmann/json (no downloads at configure time) \\
\code{examples/} & eleven ready-to-run JSON decks plus Gmsh meshes (Section~\ref{sec:examples}) \\
\code{tests/}    & unit and integration tests registered with CTest \\
\code{documents/}& this manual \\
\bottomrule
\end{tabularx}
\end{table}

%==================================================================
\section{Installation}
%==================================================================

\subsection{Requirements}

\begin{itemize}[nosep]
  \item CMake $\geq$ 3.21;
  \item a C++20 compiler (GCC, Clang or MSVC);
  \item nothing else for the base build -- Eigen and nlohmann/json
        are vendored in \code{external/}.
\end{itemize}

Optional features and their extra requirements:

\begin{itemize}[nosep]
  \item \textbf{OpenMP assembly} -- an OpenMP-capable compiler
        (enabled by default, see below);
  \item \textbf{Intel MKL PARDISO} direct solver -- an Intel oneAPI
        or standalone MKL installation;
  \item \textbf{AMGCL} iterative solver -- an AMGCL source tree and
        OpenMP;
  \item \textbf{CUDA assembly} -- the CUDA toolkit;
  \item \textbf{NVIDIA cuDSS} direct solver -- the cuDSS library
        and the CUDA toolkit.
\end{itemize}

\subsection{CMake options}

\begin{table}[htbp]
\centering
\caption{CMake configure options. ``Compile definition'' is the
preprocessor symbol the option enables in the code.}
\label{tab:cmake}
\footnotesize
\begin{tabularx}{\textwidth}{llX}
\toprule
\textbf{Option (default)} & \textbf{Compile definition} & \textbf{Effect / additional variables} \\
\midrule
\code{PARALLEL\_ASSEMBLE} (\code{ON}) & \code{PERIX\_PARALLEL\_ASSEMBLE} & builds the OpenMP assembly backend; needed for \code{"assemble": "openmp"} \\
\code{USE\_ONEAPI} (\code{OFF}) & \code{PERIX\_HAS\_PARDISO} & builds the MKL PARDISO backend; set \code{ONEAPI\_DIR} (or the \code{MKLROOT} environment variable) to the oneAPI/MKL root \\
\code{USE\_AMGCL} (\code{OFF}) & \code{PERIX\_HAS\_AMGCL} & builds the AMGCL CPU/OpenMP backend; set \code{AMGCL\_DIR} to the AMGCL source or install root; compiled with \code{AMGCL\_NO\_BOOST}; requires OpenMP \\
\code{CUDA\_ASSEMBLE} (\code{OFF}) & \code{PERIX\_CUDA\_ASSEMBLE} & builds the CUDA assembly backend; needed for \code{"assemble": "cuda"} \\
\code{USE\_CUDSS} (\code{OFF}) & \code{PERIX\_HAS\_CUDSS} & builds the NVIDIA cuDSS backend; set \code{CUDSS\_DIR} to the cuDSS root (and \code{CUDA\_DIR} if the toolkit is not in a default location) \\
\bottomrule
\end{tabularx}
\end{table}

\subsection{Building}

\begin{lstlisting}[language=bash]
git clone https://github.com/MatMechLab/PeriX.git
cd PeriX
cmake -S . -B build -DCMAKE_BUILD_TYPE=Release
cmake --build build -j
\end{lstlisting}

\noindent
The executable is \code{build/perix}; it links against the static
library \code{perix\_core} that contains the entire framework.
A fully featured configure line looks like:

\begin{lstlisting}[language=bash]
cmake -S . -B build -DCMAKE_BUILD_TYPE=Release \
  -DUSE_ONEAPI=ON -DONEAPI_DIR=/opt/intel/oneapi \
  -DUSE_AMGCL=ON  -DAMGCL_DIR=$HOME/amgcl \
  -DCUDA_ASSEMBLE=ON -DUSE_CUDSS=ON -DCUDSS_DIR=/opt/nvidia/cudss
\end{lstlisting}

\subsection{Running the test suite}

The repository registers 29 CTest cases grouped under the labels
\code{unit}, \code{integration}, \code{examples} and \code{amgcl}:

\begin{lstlisting}[language=bash]
cd build
ctest -j 16 --output-on-failure          # everything
ctest -L unit                            # only the unit tests
ctest -L examples                        # run the bundled decks
\end{lstlisting}

%==================================================================
\section{Running \PeriX{}}
%==================================================================

\subsection{Command line}

\begin{lstlisting}[language=bash]
perix -i input.json              # parse, validate and solve
perix -i input.json --read-only  # parse and validate only, no solve
\end{lstlisting}

The command line accepts exactly two options:

\begin{description}[nosep]
  \item[\code{-i <file>}] the JSON input file. The name
    \emph{must} end in \code{.json}; anything else is rejected.
  \item[\code{--read-only}] stop after reading and validating the
    input file. Useful to syntax-check a deck without running it.
\end{description}

Any parse or validation error prints a descriptive message and the
process exits with a non-zero status. On startup \PeriX{} prints a
banner with the version (0.10) and release date (2026-04-16).

\subsection{Files written by a run}

For an input file \code{job.json} (``stem'' = \code{job}):

\begin{table}[htbp]
\centering
\caption{Output files produced for input \code{job.json}.}
\label{tab:outfiles}
\footnotesize
\begin{tabularx}{\textwidth}{lX}
\toprule
\textbf{File} & \textbf{Contents} \\
\midrule
\code{job-pdmesh.vtu} & the PD point cloud (always written after mesh generation) \\
\code{job-mesh.vtu} & the background mesh; only when \code{"savemesh": true} in the \code{Mesh} block \\
\code{job-PD-results.vtu} & static analysis: nodal solution and projected fields on the PD points \\
\code{job-PD-results-00000.vtu}, \ldots & transient analysis: one file per saved step \\
\code{job-PD-results.pvd} & ParaView time-series index over the PD result files \\
\code{job-Element-results[-NNNNN].vtu} & cell-data view of the same results on the background mesh, plus \code{job-Element-results.pvd} \\
\code{job.e} & ExodusII database; only when \code{"format"} is \code{"exodus"} or \code{"both"} \\
\bottomrule
\end{tabularx}
\end{table}

Transient runs save step 0, every \code{interval}-th step, and the
final step (Section~\ref{sec:output}).

%==================================================================
\section{The Peridynamic Differential Operator}
\label{sec:pddo}
%==================================================================

This section states the PDDO construction exactly as implemented in
\code{src/PDMesh/CalcPDOperators.cpp}.

\subsection{Families, bonds and the horizon}

Every PD point $i$ owns a family
$\mathcal{H}_i=\{\,j \neq i : |\boldsymbol{\xi}_{ij}|\le\delta_i\,\}$ with
$\boldsymbol{\xi}_{ij}=\mathbf{x}_j-\mathbf{x}_i$ the \emph{bond}
vector. The horizon is
\begin{equation}
  \delta = m \cdot \Delta x,
\end{equation}
where $m$ is the input key \code{HorizonRadiusFactor} and
$\Delta x$ the characteristic grid spacing of the mesh. With
\code{VariableHorizon} enabled the per-point horizon becomes
$\delta_i=\delta\, s_i/s_{\mathrm{ref}}$
(Section~\ref{sec:varhorizon}).

\subsection{Weight and partial-volume correction}

Each bond carries a Gaussian influence weight
\begin{equation}
  w(|\boldsymbol{\xi}|)
  = \exp\!\left[-\left(\frac{2|\boldsymbol{\xi}|}{\delta}\right)^{2}\right],
  \label{eq:weight}
\end{equation}
and a Madenci-style partial-volume correction factor
\begin{equation}
  v_c(|\boldsymbol{\xi}|) =
  \begin{cases}
    1, & |\boldsymbol{\xi}| \le \delta - \Delta x/2,\\[2pt]
    \dfrac{\delta - |\boldsymbol{\xi}| + \Delta x/2}{\Delta x}
      \ \text{clamped to } [0,1], & \text{otherwise},
  \end{cases}
  \label{eq:volcorr}
\end{equation}
which linearly fades the volume of neighbours whose cells straddle
the horizon sphere. The corrected neighbour volume
$v_c\,V_j$ is used everywhere below.

\subsection{Moment matrix and operator weights}

Let $N$ be the polynomial order (input key \code{Order}, minimum 2,
maximum 6). The monomial basis on the \emph{normalised} bond
$\boldsymbol{\eta}=\boldsymbol{\xi}/\delta$ contains every
multi-index $q=(q_1,q_2,q_3)$ with $1 \le q_1+q_2+q_3 \le N$
($q_3=0$ in 2-D); the constant term is excluded because the
operator acts on differences $u_j-u_i$. For each point $i$ the
symmetric moment matrix
\begin{equation}
  A_{ab} \;=\; \sum_{j\in\mathcal{H}_i}
      w\,v_c\,V_j\;
      \boldsymbol{\eta}^{\,q_a}\,\boldsymbol{\eta}^{\,q_b}
  \label{eq:moment}
\end{equation}
is assembled, where
$\boldsymbol{\eta}^{\,q}=\eta_x^{q_1}\eta_y^{q_2}\eta_z^{q_3}$.
The coefficient matrix $a$ solves
\begin{equation}
  A\,a = B, \qquad
  B_{ab} = q_1!\,q_2!\,q_3!\;\delta_{ab},
\end{equation}
via a full-pivot LU decomposition; the solution is accepted only if
the relative residual is below $10^{-8}$, otherwise the point is
reported as degenerate and the run aborts (with one exception, see
Section~\ref{sec:ghost}). Finally the operator weight of bond
$(i,j)$ for derivative $q_k$ is
\begin{equation}
  g_i^{\,q_k}(\mathbf{x}_j) \;=\;
  \frac{w\,v_c\,V_j}{\delta^{\,|q_k|}}
  \sum_{a} a_{ak}\,\boldsymbol{\eta}^{\,q_a},
  \label{eq:gfun}
\end{equation}
and any partial derivative of a field $u$ is evaluated by the sum
\eqref{eq:pddo-idea}. Note that $V_j$ is baked into
$g_i^{\,q_k}$ -- kernels never multiply by the neighbour volume
again.

Operators are stored under human-readable names:
\code{d/dx}, \code{d/dy}, \code{d2/dx2}, \code{d2/dxdy},
\code{d2/dy2}, \code{d3/dx3}, \ldots{} -- one per multi-index up to
order $N$. Order $N=2$ (the minimum the validator accepts) already
provides all first and second derivatives, which is what every
bundled kernel consumes; higher orders increase the accuracy of
those derivatives and make higher derivatives available.

\subsection{Cracks and modified families}

When \code{MeshModify} defines pre-existing cracks
(Section~\ref{sec:meshmodify}), bonds that geometrically intersect
a crack segment are excluded from both the moment matrix
\eqref{eq:moment} and the operator sums \eqref{eq:gfun} --
\emph{unless} the treatment is \code{force\_only}, in which case
the bond stays in the family and the operators, and the fracture
kernels instead start it with zero mechanical health
(Section~\ref{sec:fracture}).

%==================================================================
\section{Meshes and the PD Point Cloud}
\label{sec:mesh}
%==================================================================

\PeriX{} first builds (or imports) a background mesh, then collapses
every cell to one PD point at its centroid carrying the cell
measure as its volume $V_i$. All physics lives on the PD cloud; the
background mesh is kept only for visualisation
(\code{*-Element-results*.vtu}).

\subsection{Built-in structured mesh (\code{"type": "perix"})}

Generates a uniform grid of \code{quad4} cells in 2-D or
\code{hex8} cells in 3-D on the box
$[x_{\min},x_{\max}]\times[y_{\min},y_{\max}]
 (\times[z_{\min},z_{\max}])$ with $n_x\times n_y(\times n_z)$
cells. PD points sit at cell centres and carry
$V_i=\Delta X\,\Delta Y(\,\Delta Z)$.

\subsection{Imported Gmsh mesh (\code{"type": "gmsh"})}

Reads an \emph{ASCII} Gmsh MSH~4.1 file (binary files are
rejected). Supported element types are \code{POINT},
\code{EDGE2}, \code{TRI3}, \code{QUAD4}, \code{TET4} and
\code{HEX8}; mixing different element types of the top dimension in
one file is rejected. The file must contain a
\code{\$PhysicalNames} section -- physical groups become \PeriX{}
node groups usable in boundary conditions. PD points sit at cell
centroids with $V_i$ the cell area/volume; the characteristic
spacing used for the horizon is the median of
$V_i^{1/\dim}$.

\subsection{Built-in disc mesh (\code{"type": "circle"})}

Generates a 2-D disc of given \code{radius} as \code{layers}
concentric rings of cells around the centre; ring thickness (and
hence local point density) is uniform in radius. Each ring's
circumference-based point count is rounded to a positive multiple of
four, and the points are offset by half an angular step. This gives
fourfold rotational symmetry without placing points directly on the
coordinate axes. The innermost ring is exposed as the four-point
group \code{anchornodes}. The outer rim is
wrapped in ghost points when \code{boundary} is \code{true} (or the
string \code{"outer"}), producing the groups \code{outernodes}
(bulk rim) and \code{outernodes\_ghost} (fictitious layer).

\subsection{Ghost boundary layer}
\label{sec:ghost}

Boundary conditions in a nonlocal method need a \emph{fictitious
layer} of points outside the physical domain. \PeriX{} adds
\code{ghost\_layer} rings (default 1) of ghost points on the outer
boundary:

\begin{itemize}[nosep]
  \item structured meshes: ghosts are placed half a cell outside
        each boundary face;
  \item imported meshes: ghosts are the mirror images of the
        near-boundary bulk points, reflected across the boundary
        face plane;
  \item disc meshes: ghosts continue the outermost ring pattern.
\end{itemize}

Each boundary wall \code{<wall>} yields two node groups:
\code{<wall>nodes\_ghost} (the ghost points) and
\code{<wall>nodes} (their mirror bulk partners). Structured walls
are named \code{left}, \code{right}, \code{bottom}, \code{top}
(plus \code{back}, \code{front} in 3-D); e.g.\
\code{leftnodes\_ghost}. Imported meshes keep their Gmsh physical
names and the same \code{\_ghost} suffix convention. The group
\code{alldomain} always addresses every bulk (non-ghost) point.

Ghost points participate in neighbour families and PDDO operators
like any other point; their rows in the global system are what the
boundary conditions overwrite (Section~\ref{sec:bc}). Points in
ghost layers deeper than the first that end up with a degenerate
moment matrix receive zero operators instead of aborting the run
(their rows are fully replaced by BCs anyway); a degenerate
\emph{bulk} point is always a fatal error.

\subsection{Variable horizon}
\label{sec:varhorizon}

With \code{"VariableHorizon": true} in \code{PDMesh}, each point
gets its own horizon
\begin{equation}
  \delta_i=\delta\,\frac{s_i}{s_{\mathrm{ref}}},\qquad
  s_i=V_i^{1/\dim},\qquad
  s_{\mathrm{ref}}=\operatorname{median}_k s_k,
\end{equation}
so the physical horizon scales with the local point spacing on
graded meshes (e.g.\ the \code{circle} mesh). A pair $(i,j)$ is
bonded if
$|\boldsymbol{\xi}_{ij}|\le\max(\delta_i,\delta_j)$, keeping bonds
symmetric. Fracture kernels then use the pairwise mean
$\delta_b=\tfrac12(\delta_i+\delta_j)$ in the critical-stretch
formula.

\subsection{Pre-existing cracks (\code{MeshModify})}
\label{sec:meshmodify}

The \code{MeshModify} block (allowed only together with a fracture
kernel) cuts cracks into the initial bond network. Cracks can be
given as explicit segments (\code{Cracks}) or generated from
parametric presets (\code{Presets}, types \code{center\_crack} and
\code{edge\_crack}); Section~\ref{sec:input-meshmodify} lists every
key. With the only public treatment, \code{force\_only}, severed
bonds keep contributing to the PDDO operators (so gradients stay
accurate near the crack faces) while the fracture kernels start
them at zero mechanical health -- the crack is invisible to
transport terms but mechanically open from step one.

%==================================================================
\section{JSON Input Reference}
\label{sec:input}
%==================================================================

\subsection{General structure and validation}

The input is a single JSON object. Twelve top-level blocks are
recognised; any other key is rejected. Validation is strict and
happens before any computation: unknown keys \emph{inside} a block,
wrong JSON types, and out-of-range values all abort the run with a
message naming the offending block and key.

\begin{table}[htbp]
\centering
\caption{Top-level blocks.}
\label{tab:blocks}
\footnotesize
\begin{tabularx}{\textwidth}{llX}
\toprule
\textbf{Block} & \textbf{Presence} & \textbf{Purpose} \\
\midrule
\code{Mesh}            & required & background mesh (Section~\ref{sec:input-mesh}) \\
\code{PDMesh}          & required & horizon, PDDO order, ghost layers (Section~\ref{sec:input-pdmesh}) \\
\code{MeshModify}      & optional & pre-existing cracks (Section~\ref{sec:input-meshmodify}) \\
\code{DOFs}            & required & symbolic field names (Section~\ref{sec:input-dofs}) \\
\code{ElmtSystem}      & required & the physics kernel and its parameters (Section~\ref{sec:elmt}) \\
\code{BCSystem}        & required & boundary conditions (Section~\ref{sec:bc}) \\
\code{ICSystem}        & optional & initial conditions (Section~\ref{sec:ic}) \\
\code{JobSystem}       & optional & static/transient switch, assembly backend (Section~\ref{sec:job}) \\
\code{TimeStepping}    & see text & time step control; required for transient jobs, forbidden for static ones \\
\code{NonlinearSolver} & optional & Newton iteration control (Section~\ref{sec:newton}) \\
\code{LinearSolver}    & optional & linear solver selection (Section~\ref{sec:linsolver}) \\
\code{OutputSystem}    & optional & output cadence, fields, format (Section~\ref{sec:output}) \\
\bottomrule
\end{tabularx}
\end{table}

With \code{--read-only} (or when only the mesh is being inspected)
just the \code{Mesh} block is mandatory; a normal solve requires
\code{Mesh}, \code{PDMesh}, \code{DOFs}, \code{ElmtSystem} and
\code{BCSystem}.

\subsection{\code{Mesh}}
\label{sec:input-mesh}

The mandatory key \code{type} selects one of three generators:
\code{"perix"} (structured box), \code{"gmsh"} (import), or
\code{"circle"} (2-D disc).

\subsubsection*{\code{"type": "perix"} -- structured box}

\begin{table}[htbp]
\centering
\caption{\code{Mesh} keys for \code{type=perix}.}
\footnotesize
\begin{tabularx}{\textwidth}{lllX}
\toprule
\textbf{Key} & \textbf{Type} & \textbf{Default} & \textbf{Meaning / constraints} \\
\midrule
\code{dim}      & int    & required & 2 or 3 \\
\code{meshtype} & string & required & \code{"quad4"} in 2-D, \code{"hex8"} in 3-D (enforced) \\
\code{nx},\code{ny} & int & required & cells per direction, $\geq 1$ \\
\code{nz}       & int    & required in 3-D & rejected in 2-D, as are \code{zmin}/\code{zmax} \\
\code{xmin},\code{xmax} & number & 0, 1 & domain extent; every minimum must be $<$ its maximum \\
\code{ymin},\code{ymax} & number & 0, 1 & \\
\code{zmin},\code{zmax} & number & 0, 1 & 3-D only \\
\code{savemesh} & bool   & \code{false} & write \code{<stem>-mesh.vtu} \\
\bottomrule
\end{tabularx}
\end{table}

\subsubsection*{\code{"type": "gmsh"} -- imported mesh}

\begin{table}[htbp]
\centering
\caption{\code{Mesh} keys for \code{type=gmsh}.}
\footnotesize
\begin{tabularx}{\textwidth}{lllX}
\toprule
\textbf{Key} & \textbf{Type} & \textbf{Default} & \textbf{Meaning / constraints} \\
\midrule
\code{file}     & string & required & path to an \emph{ASCII} Gmsh MSH 4.1 file; binary files are rejected \\
\code{savemesh} & bool   & \code{false} & write \code{<stem>-mesh.vtu} \\
\bottomrule
\end{tabularx}
\end{table}

The file must contain \code{\$PhysicalNames}; the physical groups
become the node groups referenced by boundary conditions. Supported
cell types: \code{POINT}, \code{EDGE2}, \code{TRI3}, \code{QUAD4},
\code{TET4}, \code{HEX8}; mixing different cell types of the
highest dimension is rejected.

\subsubsection*{\code{"type": "circle"} -- 2-D disc}

\begin{table}[htbp]
\centering
\caption{\code{Mesh} keys for \code{type=circle}.}
\footnotesize
\begin{tabularx}{\textwidth}{lllX}
\toprule
\textbf{Key} & \textbf{Type} & \textbf{Default} & \textbf{Meaning / constraints} \\
\midrule
\code{radius}   & number & required & disc radius, $>0$ \\
\code{layers}   & int    & required & number of concentric point rings, $\geq 1$ \\
\code{center}   & array  & origin   & 2 or 3 numbers \\
\code{boundary} & bool or string & \code{true} & \code{true}: wrap the rim in ghost points with group base name \code{outer}; a non-empty string sets a custom base name; \code{false}: no ghost rim \\
\code{savemesh} & bool   & \code{false} & write \code{<stem>-mesh.vtu} \\
\bottomrule
\end{tabularx}
\end{table}

The rim groups are \code{<base>nodes} (outermost bulk ring) and
\code{<base>nodes\_ghost} (fictitious rim), i.e.\
\code{outernodes} / \code{outernodes\_ghost} by default.

\subsection{\code{PDMesh}}
\label{sec:input-pdmesh}

\begin{table}[htbp]
\centering
\caption{\code{PDMesh} keys.}
\footnotesize
\begin{tabularx}{\textwidth}{lllX}
\toprule
\textbf{Key} & \textbf{Type} & \textbf{Default} & \textbf{Meaning / constraints} \\
\midrule
\code{HorizonRadiusFactor} & number & required & $m$ in $\delta=m\,\Delta x$; must be positive \\
\code{Order} & int & required & PDDO polynomial order $N$; the validator accepts $2 \le N \le 6$ (the published kernels require at least 2) \\
\code{VariableHorizon} & bool & \code{false} & per-point horizon scaling for graded meshes (Section~\ref{sec:varhorizon}) \\
\code{ghost\_layer} & int & 1 & number of fictitious boundary layers, $\geq 1$ \\
\bottomrule
\end{tabularx}
\end{table}

A useful rule of thumb is $m \approx 3$ with \code{Order} 2; the
bundled examples use horizon factors between 2.03 and 3.03.

\subsection{\code{MeshModify}}
\label{sec:input-meshmodify}

Allowed only when the element is one of the fracture kernels.

\begin{table}[htbp]
\centering
\caption{\code{MeshModify} keys.}
\footnotesize
\begin{tabularx}{\textwidth}{lllX}
\toprule
\textbf{Key} & \textbf{Type} & \textbf{Default} & \textbf{Meaning / constraints} \\
\midrule
\code{type} & string & required & must be \code{"pre\_existing\_cracks"} \\
\code{treatment} & string & required & must be \code{"force\_only"}: severed bonds stay in the PDDO operators but start with zero mechanical health \\
\code{Cracks} & array & optional & explicit 2-D crack segments, see below \\
\code{Presets} & array & optional & parametric crack generators, see below \\
\bottomrule
\end{tabularx}
\end{table}

At least one of \code{Cracks} / \code{Presets} must be present.
Each \code{Cracks} item is an object
\code{\{x1, y1, x2, y2\}} (numbers, required; the two endpoints
must differ) plus an optional \code{label} string.

Each \code{Presets} item has \code{type} \code{"center\_crack"} or
\code{"edge\_crack"}:

\begin{table}[htbp]
\centering
\caption{\code{Presets} item keys.}
\footnotesize
\begin{tabularx}{\textwidth}{llX}
\toprule
\textbf{Key} & \textbf{Applies to} & \textbf{Meaning / constraints} \\
\midrule
\code{center} or \code{center\_x}/\code{center\_y}/\code{center\_z} & center\_crack & crack midpoint; the array form and the per-component form are mutually exclusive; components default to the domain centre \\
\code{length} & both & crack length, $>0$; required \\
\code{width} & both & out-of-plane extrusion width (3-D), $>0$, optional \\
\code{angle\_degrees} or \code{angle\_radians} & both & in-plane inclination; mutually exclusive; \code{edge\_crack} defaults to the inward normal of its side ($0$ for \code{left}, $\pi$ for \code{right}, $\pi/2$ for \code{bottom}, $-\pi/2$ for \code{top}) \\
\code{normal}, \code{axis} & center\_crack & optional 3-vectors selecting the crack-plane normal / in-plane axis in 3-D; without them (and without \code{width}) a 3-D centre crack is extruded through the whole thickness \\
\code{side} & edge\_crack & required: \code{"left"}, \code{"right"}, \code{"bottom"} or \code{"top"} \\
\code{position} & edge\_crack & coordinate of the mouth along the chosen side \\
\code{label} & both & optional name echoed in the log \\
\bottomrule
\end{tabularx}
\end{table}

\subsection{\code{DOFs}}
\label{sec:input-dofs}

A non-empty array of unique, non-empty field-name strings, e.g.\
\code{["c","mu","ux","uy"]}. The names are \emph{not} free: every
kernel prescribes its canonical list and order
(Table~\ref{tab:kernels}); the validator rejects any deviation.
Boundary and initial conditions refer to these names in their
\code{dofs}/\code{dof} keys.

\subsection{Cross-block rules}
\label{sec:crossrules}

After all blocks are validated individually, the following global
rules are enforced (each quotes the actual error condition):

\begin{itemize}[nosep]
  \item a \code{transient} job requires a \code{TimeStepping}
        block; a \code{static} job must not have one;
  \item \code{explicit\_pddo\_frac} requires a transient job;
  \item \code{explicit\_pddo\_frac} is matrix-free and rejects both
        \code{LinearSolver} and \code{NonlinearSolver} blocks, and
        also \code{pdtraction} boundary conditions (they need an
        assembled system);
  \item \code{"field": "velocity"} initial conditions are accepted
        only with the explicit central-difference element;
  \item \code{speciesflux} boundary conditions require a species
        transport kernel (\code{cahnhilliard},
        \code{stress\_cahnhilliard} or
        \code{frac\_stress\_cahnhilliard});
  \item \code{MeshModify} requires a fracture kernel
        (\code{pddo\_dynamic\_frac}, \code{explicit\_pddo\_frac} or
        \code{frac\_stress\_cahnhilliard}).
\end{itemize}

%==================================================================
\section{Physics Kernels}
\label{sec:elmt}
%==================================================================

\subsection{The \code{ElmtSystem} block}

\code{ElmtSystem} is an object containing \emph{exactly one} entry.
The entry's key is a free label (it only appears in error
messages); its value is an object with the two keys \code{type}
(one of the seven kernel names) and \code{params} (an object,
always required):

\begin{lstlisting}[language=json]
"ElmtSystem": {
  "myblock": {
    "type": "poisson",
    "params": { "sigma": 1.0, "f": 1.0 }
  }
}
\end{lstlisting}

The top-level \code{DOFs} array must match the kernel's canonical
field list of Table~\ref{tab:kernels} exactly -- same names, same
order. For the mechanical kernels on a mesh whose dimension cannot
be determined at validation time (an imported Gmsh file), both the
2-D and the 3-D DoF lists are accepted and the actual mesh decides.

Every kernel exposes \emph{projections}: derived nodal fields
(gradients, stresses, damage, \ldots) computed after each converged
step and written to the output files when requested by name in
\code{OutputSystem.Fields} (Section~\ref{sec:output}).

\subsection{\code{poisson} -- steady Poisson equation}

Solves
\begin{equation}
  \sigma \Lap u + f = 0
\end{equation}
on the PD cloud, with the Laplacian evaluated through the PDDO
second-derivative operators
($\Lap = \code{d2/dx2}+\code{d2/dy2}(+\code{d2/dz2})$). The kernel
is linear; a static job converges in one Newton iteration.

\begin{table}[htbp]
\centering
\caption{\code{poisson} parameters.}
\footnotesize
\begin{tabularx}{\textwidth}{lllX}
\toprule
\textbf{Key} & \textbf{Type} & \textbf{Constraint} & \textbf{Meaning} \\
\midrule
\code{sigma} & number & $>0$, required & conductivity $\sigma$ \\
\code{f}     & number & required       & constant source term $f$ \\
\bottomrule
\end{tabularx}
\end{table}

\noindent\textbf{Projections:} \code{gradient} $=\Grad u$ and
\code{flux} $=-\sigma\Grad u$ (vector fields).

\subsection{\code{diffusion} -- transient diffusion}

Solves
\begin{equation}
  \frac{\partial c}{\partial t} = D \Lap c + f
\end{equation}
with a backward-Euler residual
$R_i=(c_i-c_i^{\text{old}})/\Delta t - D\Lap c_i - f$. The kernel
is transient-only: it refuses to run without a positive time step,
so the job must be \code{transient} with a \code{TimeStepping}
block.

\begin{table}[htbp]
\centering
\caption{\code{diffusion} parameters.}
\footnotesize
\begin{tabularx}{\textwidth}{lllX}
\toprule
\textbf{Key} & \textbf{Type} & \textbf{Constraint} & \textbf{Meaning} \\
\midrule
\code{D} & number & $>0$, required & diffusivity \\
\code{f} & number & required       & constant source term \\
\bottomrule
\end{tabularx}
\end{table}

\noindent\textbf{Projections:} \code{gradient} $=\Grad c$,
\code{flux} $=-D\Grad c$.

\subsection{\code{cahnhilliard} -- phase separation}
\label{sec:ch}

Solves the Cahn--Hilliard equation in the mixed two-field form
\begin{align}
  \frac{\partial c}{\partial t} &= \Div\!\bigl[M(c)\,\Grad\mu\bigr],\\
  \mu &= f'(c) - \kappa \Lap c,
\end{align}
with the regular-solution free energy
\begin{equation}
  f'(c)=\ln\frac{c}{1-c}+\chi(1-2c),\qquad
  f''(c)=\frac{1}{c(1-c)}-2\chi ,
\end{equation}
and the degenerate mobility $M(c)=M_0\,c(1-c)$. The concentration
is clamped to $[\varepsilon,\,1-\varepsilon]$ with
$\varepsilon=10^{-8}$ inside the logarithm and the mobility, so the
scheme survives $c\to 0,1$. The mobility is Picard-frozen: it is
evaluated at the current iterate and treated as constant in the
Jacobian.

Three implementation details matter for accuracy and mass
conservation:
\begin{itemize}[nosep]
  \item \textbf{Conservative flux.} The transport term is
    discretised with the antisymmetrised bond weight
    \begin{equation}
      w_{ij}=\tfrac12\Bigl[M_i\,g^{2}_{ij}
        +\tfrac{V_j}{V_i}\,M_j\,g^{2}_{ji}\Bigr],\qquad
      \Div(M\Grad\mu)\big|_i \approx \sum_j w_{ij}(\mu_j-\mu_i),
    \end{equation}
    where $g^{2}_{ij}$ is the PDDO Laplacian weight of bond
    $(i,j)$. This makes the exchange between $i$ and $j$ exactly
    antisymmetric after multiplication by $V_i$, so the total mass
    $\sum_i V_i c_i$ is conserved to solver tolerance.
  \item \textbf{No flux through ghosts.} Bonds that reach into the
    fictitious boundary layer are dropped from the species flux and
    from the $\kappa\Lap c$ term, closing the domain (zero-flux
    walls) regardless of the boundary conditions applied to the
    ghost rows.
  \item \textbf{Laplacian calibration.} Each node's Laplacian
    weights are rescaled by
    $\lambda_i = 2\dim / \sum_j g^{2}_{ij}\,
     |\mathbf{x}_j-\mathbf{x}_i|^{2}$
    (the exact Laplacian of $|\mathbf{x}|^{2}$ is $2\dim$), which
    repairs the one-sided stencils next to the boundary.
\end{itemize}

\begin{table}[htbp]
\centering
\caption{\code{cahnhilliard} parameters.}
\footnotesize
\begin{tabularx}{\textwidth}{lllX}
\toprule
\textbf{Key} & \textbf{Type} & \textbf{Constraint} & \textbf{Meaning} \\
\midrule
\code{chi}   & number & required        & Flory interaction parameter $\chi$ \\
\code{kappa} & number & $>0$, required  & gradient-energy coefficient $\kappa$ \\
\code{M}     & number & $>0$, required  & mobility scale $M_0$ in $M(c)=M_0\,c(1-c)$ \\
\bottomrule
\end{tabularx}
\end{table}

\noindent\textbf{Projections:} \code{fPrime} $=f'(c)$ and
\code{fDoublePrime} $=f''(c)$ (scalars).

\subsection{Bond-based elasticity and fracture: common ground}
\label{sec:fracture}

The four mechanical kernels share one constitutive core. With
displacement vector $\mathbf{u}$ and the isotropic stiffness
constants
\begin{equation}
  \begin{aligned}
  \text{plane stress:}\quad &C_{11}=\frac{E}{1-\nu^{2}},\quad
    C_{12}=\nu C_{11},\quad C_{66}=\frac{E}{2(1+\nu)},\\
  \text{plane strain / 3-D:}\quad &C_{11}=\lambda+2\mu_s,\quad
    C_{12}=\lambda,\quad C_{66}=\mu_s,
  \end{aligned}
\end{equation}
($\lambda,\mu_s$ the Lam\'e constants), the discrete Navier
operator at point $i$ is assembled bond-wise from the PDDO
second-derivative weights $G^{xx}_{ij},G^{yy}_{ij},G^{xy}_{ij},
\ldots$ In 2-D:
\begin{equation}
  \begin{aligned}
  L_x\big|_i &= \sum_j \mu_{ij}\Bigl[
     (C_{11}G^{xx}_{ij}+C_{66}G^{yy}_{ij})\,\Delta u_x
     +(C_{12}+C_{66})\,G^{xy}_{ij}\,\Delta u_y \Bigr],\\
  L_y\big|_i &= \sum_j \mu_{ij}\Bigl[
     (C_{12}+C_{66})\,G^{xy}_{ij}\,\Delta u_x
     +(C_{66}G^{xx}_{ij}+C_{11}G^{yy}_{ij})\,\Delta u_y \Bigr],
  \end{aligned}
  \label{eq:navier}
\end{equation}
with $\Delta u_\bullet = u_{\bullet j}-u_{\bullet i}$ and the same
pattern extended in 3-D. Here $\mu_{ij}\in\{0,1\}$ is the
\emph{bond health}: 1 while intact, 0 once broken.

\paragraph{Critical stretch.}
A bond breaks when its stretch
$s=(|\boldsymbol{\xi}+\boldsymbol{\eta}_u|-|\boldsymbol{\xi}|)/|\boldsymbol{\xi}|$
(with $\boldsymbol{\eta}_u$ the relative displacement) exceeds the
critical value calibrated from the energy release rate $G_c$:
\begin{equation}
  s_{cr}=
  \begin{cases}
    \sqrt{\dfrac{4\pi G_c}{9 E \delta}} & \text{plane stress},\\[8pt]
    \sqrt{\dfrac{5\pi G_c}{12 E \delta}} & \text{plane strain},\\[8pt]
    \sqrt{\dfrac{5 G_c}{6 E \delta}} & \text{3-D},
  \end{cases}
  \label{eq:scr}
\end{equation}
With \code{"tension\_only": true} only tensile stretches
($s>s_{cr}$) break bonds; with \code{false} compressive bonds break
too ($|s|>s_{cr}$). Bond failure is irreversible. With variable
horizons the bond uses $\delta_b=\tfrac12(\delta_i+\delta_j)$ in
\eqref{eq:scr}. The nodal \code{damage} projection is
\begin{equation}
  d_i = 1-\frac{\text{intact bonds of } i}{\text{family size of } i}
  \in [0,1].
\end{equation}

\subsection{\code{pddo\_dynamic\_frac} -- implicit dynamic fracture}
\label{sec:pddofrac}

Solves $\rho\,\ddot{\mathbf u} = \mathbf{L}(\mathbf u)$ (no body
force in the public kernel) with an implicit backward-Euler
integrator embedded in the residual:
\begin{equation}
  \mathbf R_i = \mathbf L(\mathbf u)\big|_i
  - \rho\,
  \frac{\mathbf u_i-\mathbf u_i^{n-1}-\mathbf v_i^{n-1}\Delta t}
       {\Delta t^{2}} ,
\end{equation}
which Newton drives to zero each step; on convergence the velocity
is committed as
$\mathbf v^{n}=(\mathbf u^{n}-\mathbf u^{n-1})/\Delta t$. Bond
failure is evaluated once per \emph{converged} step on the
committed displacement (never inside the Newton loop), so the
tangent stays consistent within a step.

\begin{table}[htbp]
\centering
\caption{\code{pddo\_dynamic\_frac} / \code{explicit\_pddo\_frac} parameters.}
\label{tab:fracparams}
\footnotesize
\begin{tabularx}{\textwidth}{lllX}
\toprule
\textbf{Key} & \textbf{Type} & \textbf{Constraint} & \textbf{Meaning} \\
\midrule
\code{E}   & number & $>0$, required & Young's modulus \\
\code{nu}  & number & $\in(-1,0.5)$, required & Poisson's ratio \\
\code{rho} & number & $>0$, required & mass density \\
\code{Gc} / \code{G0} & number & $>0$, exactly one & critical energy release rate (alias pair) \\
\code{state} & string & required in 2-D, rejected in 3-D & \code{"plane\_stress"} or \code{"plane\_strain"} \\
\code{tension\_only} & bool & required & break bonds only in tension \\
\code{damage} & bool & required & master switch for bond breaking \\
\bottomrule
\end{tabularx}
\end{table}

\noindent\textbf{Projections:} \code{strain}, \code{stress}
(symmetric tensors, computed from PDDO displacement gradients and
the elastic law), \code{vonMises} (scalar), \code{damage}
(scalar).

\subsection{\code{explicit\_pddo\_frac} -- explicit dynamic fracture}
\label{sec:explicit}

The same physics as Section~\ref{sec:pddofrac} integrated with the
central-difference scheme
\begin{equation}
  \mathbf u^{n+1} = 2\mathbf u^{n}-\mathbf u^{n-1}
  + \Delta t^{2}\,
  \frac{\mathbf L(\mathbf u^{n})}{\rho}.
\end{equation}
No linear system is ever assembled -- the input must not contain
\code{LinearSolver} or \code{NonlinearSolver} blocks, and
\code{pdtraction} boundary conditions are rejected
(Section~\ref{sec:crossrules}). Parameters are identical to
Table~\ref{tab:fracparams}. Additional characteristics:

\begin{itemize}[nosep]
  \item \textbf{Stability.} The step must satisfy a CFL-type bound;
        the kernel estimates
        $\Delta t \lesssim \Delta x/\sqrt{C_{11}/\rho}$ from the
        grid spacing and warns when exceeded.
  \item \textbf{Velocity ICs.} \code{"field": "velocity"} initial
        conditions are supported (only by this kernel) and seed the
        $\mathbf u^{-1}$ history so the first step moves with the
        prescribed velocity.
  \item \textbf{Damage update.} Bond breaking is evaluated every
        step on the current displacement.
  \item \textbf{Performance.} Per-bond operator weights are cached
        in flat compressed arrays at start-up (with an overflow
        guard for clouds beyond $2^{31}$ bonds). When compiled with
        \code{CUDA\_ASSEMBLE} and \code{"assemble": "cuda"}, the
        bond loop runs on the GPU; if the GPU path fails it falls
        back to the CPU \emph{permanently} for the rest of the run,
        so the irreversible damage state never diverges between the
        two code paths.
\end{itemize}

\noindent\textbf{Projections:} \code{damage}, \code{vonMises},
\code{stress}, \code{strain}.

\subsection{\code{stress\_cahnhilliard} -- chemo-mechanics}
\label{sec:sch}

Couples Cahn--Hilliard transport (Section~\ref{sec:ch}) to linear
elasticity through the composition eigenstrain
$\boldsymbol\varepsilon^{\mathrm{ch}}=\frac{\Omega}{3}(c-c_{\mathrm{ref}})\mathbf I$
and the hydrostatic-stress term in the chemical potential:
\begin{align}
  \frac{\partial c}{\partial t}
    &= \Div\!\bigl[M(c)\,\Grad\mu\bigr],
    \qquad M(c)=D\,c(1-c),\\
  \mu &= \ln\frac{c}{1-c}+\chi(1-2c)
        -\Omega\,\sigma_h - \kappa\Lap c,\\
  \boldsymbol\sigma &= \mathbb C :
    \bigl(\boldsymbol\varepsilon-\boldsymbol\varepsilon^{\mathrm{ch}}\bigr),
  \qquad
  \Div\boldsymbol\sigma = \rho\,\ddot{\mathbf u}
  \;\;(\text{or } \mathbf 0 \text{ if } \rho = 0).
\end{align}
The hydrostatic stress is evaluated in closed form,
\begin{equation}
  \sigma_h = K_h\,\operatorname{tr}\boldsymbol\varepsilon
           + C_h\,(c-c_{\mathrm{ref}}),
\end{equation}
with coefficients that depend on the plane reduction:
\begin{equation}
  \begin{aligned}
  \text{plane stress:}\quad
    & A=\frac{E\,\Omega}{3(1-\nu)}, &
    K_h&=\frac{C_{11}+C_{12}}{3}, &
    C_h&=-\frac{2A}{3},\\
  \text{plane strain / 3-D:}\quad
    & A=\frac{E\,\Omega}{3(1-2\nu)}, &
    K_h&=\frac{C_{11}+2C_{12}}{3}=K, &
    C_h&=-A,
  \end{aligned}
  \label{eq:sigmah}
\end{equation}
where $A$ is also the eigenstress modulus that enters the momentum
residual as the body-force-like term $-A\,\Grad c$. The species
transport uses the same conservative antisymmetric flux, ghost-bond
exclusion and Laplacian calibration as the plain Cahn--Hilliard
kernel. With $\rho>0$ a backward-Euler inertia term is added to the
momentum residual exactly as in Section~\ref{sec:pddofrac}; with
$\rho=0$ (or omitted) the mechanics are quasi-static and the
displacements must be pinned somewhere to remove rigid-body modes.

\begin{table}[htbp]
\centering
\caption{\code{stress\_cahnhilliard} parameters.}
\footnotesize
\begin{tabularx}{\textwidth}{lllX}
\toprule
\textbf{Key} & \textbf{Type} & \textbf{Constraint} & \textbf{Meaning} \\
\midrule
\code{E}     & number & $>0$, required & Young's modulus \\
\code{nu}    & number & $\in(-1,0.5)$, required & Poisson's ratio \\
\code{D}     & number & $>0$, required & mobility scale in $M(c)=D\,c(1-c)$ \\
\code{Omega} & number & required & partial molar volume $\Omega$ \\
\code{cref}  & number & required & stress-free reference concentration \\
\code{chi}   & number & required & Flory interaction parameter \\
\code{kappa} & number & $>0$, required & gradient-energy coefficient \\
\code{rho}   & number & $\geq 0$, optional (default 0) & density; $>0$ adds backward-Euler inertia \\
\code{state} & string & required in 2-D, rejected in 3-D & \code{"plane\_stress"} or \code{"plane\_strain"} \\
\bottomrule
\end{tabularx}
\end{table}

\noindent\textbf{Projections:} \code{strain}, \code{stress},
\code{vonMises}, \code{sigmaH} ($\sigma_h$), \code{fPrime},
\code{fDoublePrime}, \code{gradC} ($\Grad c$).

\subsection{\code{frac\_stress\_cahnhilliard} -- chemo-mechanics with fracture}
\label{sec:fsch}

Extends Section~\ref{sec:sch} with critical-stretch bond breaking.
The design principle: \emph{damage degrades mechanics only}. Bond
health $\mu_{ij}$ multiplies

\begin{itemize}[nosep]
  \item the elastic stiffness contributions \eqref{eq:navier},
  \item the eigenstress body force $-A\,\Grad c$, and
  \item the $\Omega K_h \Div\mathbf u$ mechanical-coupling term in
        the $\mu$-equation,
\end{itemize}

while species diffusion and the $\kappa$ interface term are
\emph{never} degraded -- a cracked region still transports species.
Two fracture-specific refinements:

\begin{itemize}[nosep]
  \item \textbf{Swelling-corrected stretch.} The failure criterion
        uses the elastic part of the bond stretch,
        \begin{equation}
          s_{\mathrm{el}} = s - \frac{\Omega}{3}
          \bigl(\bar c_{ij}-c_{\mathrm{ref}}\bigr),\qquad
          \bar c_{ij}=\tfrac12(c_i+c_j),
        \end{equation}
        so homogeneous swelling does not spuriously break bonds.
  \item \textbf{Residual stiffness.} A node whose bonds have all
        failed would leave a singular block in the Jacobian; the
        stiffness gate is therefore
        $g_{\mathrm{eff}}=\max(\mu_{ij},\,k_{\mathrm{res}})$
        applied through the \code{residual\_stiffness} floor
        $k_{\mathrm{res}}\in(0,1]$ (default $10^{-8}$).
\end{itemize}

\begin{table}[htbp]
\centering
\caption{\code{frac\_stress\_cahnhilliard} parameters (in addition
to every key of \code{stress\_cahnhilliard}; here \code{rho} is
\emph{required}).}
\footnotesize
\begin{tabularx}{\textwidth}{lllX}
\toprule
\textbf{Key} & \textbf{Type} & \textbf{Constraint} & \textbf{Meaning} \\
\midrule
\code{rho}  & number & $\geq 0$, \emph{required} & 0: quasi-static (pin rigid modes!); $>0$: backward-Euler inertia \\
\code{Gc} / \code{G0} & number & $>0$, exactly one & critical energy release rate (alias pair) \\
\code{tension\_only} & bool & required & break bonds only in tension \\
\code{damage\_on} / \code{damage} & bool & exactly one & master switch for bond breaking (alias pair) \\
\code{residual\_stiffness} & number & $\in(0,1]$, default $10^{-8}$ & stiffness floor for fully severed nodes \\
\bottomrule
\end{tabularx}
\end{table}

\noindent\textbf{Projections:} those of
\code{stress\_cahnhilliard} plus \code{damage}.

%==================================================================
\section{Boundary Conditions}
\label{sec:bc}
%==================================================================

\subsection{The \code{BCSystem} block}

\code{BCSystem} is a non-empty object; every entry is one boundary
condition, named by a free label. Each entry requires \code{type}
(one of \code{dirichlet}, \code{neumann}, \code{speciesflux},
\code{pdtraction}, \code{sdtraction}) and \code{phygroup} -- the
node group it acts on (Section~\ref{sec:ghost}). In a nonlocal
method a ``boundary condition'' is a recipe for filling the rows of
the fictitious ghost points, so \code{phygroup} is usually a
\code{*nodes\_ghost} group.

Two run-time safety nets complement the schema: a Dirichlet
condition that pins a Cahn--Hilliard species DoF (\code{c} or
\code{mu}) on boundary ghosts is rejected (the species flux drops
ghost bonds, so such a pin is inert and would silently seal the
domain), and any ghost wall left without \emph{any} boundary
condition is reported at start-up, because unconstrained one-sided
PD rows are the most common cause of solver failure on refined
meshes.

\subsection{\code{dirichlet} -- essential conditions}

\begin{table}[htbp]
\centering
\caption{\code{dirichlet} keys.}
\footnotesize
\begin{tabularx}{\textwidth}{lllX}
\toprule
\textbf{Key} & \textbf{Type} & \textbf{Default} & \textbf{Meaning} \\
\midrule
\code{phygroup} & string & required & target node group \\
\code{dofs}  & array & required & non-empty list of declared DoF names \\
\code{value} & number or array & required & prescribed value $g_0$; an array must match \code{dofs} in length \\
\code{velocity} & number or array & 0 & ramp rate: $g(t)=g_0+\dot g\,t$ \\
\code{box} & object & none & axis-aligned filter \code{\{xmin,xmax,ymin,ymax,zmin,zmax\}}; only nodes inside all given bounds are affected (a selection that matches no node warns once) \\
\code{direct} & bool & \code{false} & force direct pinning (see below) \\
\bottomrule
\end{tabularx}
\end{table}

On a ghost group the default is the \emph{mirror (reflection)}
form: for each ghost $g$ with mirror bulk partner $b$ the row
enforces
\begin{equation}
  u_g + u_b = 2\,g(t)
  \qquad\Longleftrightarrow\qquad
  u\big|_{\text{wall}} = g(t),
\end{equation}
i.e.\ the wall value midway between ghost and bulk equals the
prescribed value exactly, second-order accurately. With
\code{"direct": true}, on non-ghost groups, or for ghosts without a
mirror partner, the row degenerates to the direct pin
$u_g=g(t)$. Row replacement scales the unit diagonal to the mean
magnitude of the matrix diagonal, preserving the conditioning of
the assembled system.

\subsection{\code{neumann} -- homogeneous natural conditions}

Keys: \code{phygroup}, \code{dofs}, \code{value} -- and the value
\emph{must be zero} (scalar 0 or an all-zero array); the public
generic Neumann condition is homogeneous. It enforces the
zero-normal-gradient mirror $u_g=u_b$. For inhomogeneous loads use
\code{pdtraction}/\code{sdtraction} (mechanics) or
\code{speciesflux} (species transport).

\subsection{\code{speciesflux} -- prescribed species influx}

Keys: \code{phygroup}, \code{dofs} (must be exactly
\code{["c"]}), \code{value} (the flux $j$). Requires a species
kernel. The implementation is a \emph{source form}: the ghost row
keeps the zero-gradient mirror $c_g=c_b$ (so the PDDO family of the
wall sees no artificial gradient), and the flux enters the
mass balance of the near-wall \emph{bulk} row as a source
\begin{equation}
  R_b \mathrel{-}= \frac{j}{h},
\end{equation}
where $h$ is the wall-layer thickness (cell volume divided by face
area where available, otherwise the ghost--bulk distance).
Positive \code{value} means influx: the total species content
grows at exactly $j$ per unit wall area and time.

\subsection{\code{pdtraction} and \code{sdtraction} -- strong-form tractions}

These conditions replace the ghost rows with the strong-form
traction statement
\begin{equation}
  \boldsymbol\sigma(\mathbf x_w)\cdot\mathbf n = \bar{\mathbf t},
  \qquad
  \boldsymbol\sigma=\mathbb C:\boldsymbol\varepsilon
  \;-\;
  \underbrace{A\,(c-c_{\mathrm{ref}})\,\mathbf I}_{\text{\code{sdtraction} only}},
\end{equation}
built from the PDDO first-derivative operators of the ghost point;
$\mathbf n$ is the dominant-axis outward normal of the wall. They
require an assembled linear system and are therefore invalid with
the explicit kernel. \code{pdtraction} is the purely mechanical
form; \code{sdtraction} additionally subtracts the swelling
eigenstress, which is the correct traction-free condition for the
chemo-mechanical kernels.

\begin{table}[htbp]
\centering
\caption{\code{pdtraction} / \code{sdtraction} keys.}
\footnotesize
\begin{tabularx}{\textwidth}{lllX}
\toprule
\textbf{Key} & \textbf{Type} & \textbf{pdtraction} & \textbf{sdtraction} \\
\midrule
\code{phygroup} & string & required & required \\
\code{E}, \code{nu} & numbers & required ($E>0$, $\nu\in(-1,0.5)$) & same \\
\code{state} & string & required in 2-D & optional in 2-D (default \code{plane\_stress}); rejected in 3-D for both \\
\code{value} & array[dim] & required & optional (default zero traction) \\
\code{dofs} & array & \multicolumn{2}{l}{optional, but \emph{required} for the coupled kernels; must be \code{["ux","uy"(,"uz")]}}\\
\code{Omega} & number & \multicolumn{2}{l}{chemical coupling: \code{pdtraction} needs \code{Omega}+\code{cref}+\code{c\_dof} together;}\\
\code{cref} & number & \multicolumn{2}{l}{\quad \code{sdtraction} needs \code{Omega} only (\code{cref} defaults to 0, \code{c\_dof} to \code{"c"})}\\
\code{c\_dof} & string & \multicolumn{2}{l}{must name the declared concentration field \code{"c"}}\\
\bottomrule
\end{tabularx}
\end{table}

The elastic constants of the traction condition are given
explicitly (\code{E}, \code{nu}, \code{state}) so the boundary
operator is self-contained; in practice they repeat the element
parameters.

%==================================================================
\section{Initial Conditions}
\label{sec:ic}
%==================================================================

\code{ICSystem} is a non-empty object of named entries. Every entry
requires \code{type} and targets either one DoF (\code{"dof":
"c"}) or several (\code{"dofs": [...]}) -- exactly one of the two
spellings. Common optional keys:

\begin{itemize}[nosep]
  \item \code{field}: \code{"solution"} (default) or
        \code{"velocity"}; velocity ICs are accepted only with
        \code{explicit\_pddo\_frac};
  \item \code{phygroup}: restrict to a node group (default: all
        bulk nodes).
\end{itemize}

Entries are applied in declaration order; later entries overwrite
earlier ones on overlapping nodes. Keys of kind ``scalar or
array'' accept one number (applied to every listed DoF) or an
array matching \code{dofs} in length.

\begin{table}[htbp]
\centering
\caption{IC generators. ``s/a'' = scalar or per-DoF array.}
\label{tab:ics}
\footnotesize
\begin{tabularx}{\textwidth}{l>{\raggedright\arraybackslash}X>{\raggedright\arraybackslash}X}
\toprule
\textbf{\code{type}} & \textbf{Keys} & \textbf{Field value at $\mathbf x$} \\
\midrule
\code{constant} & \code{value} (s/a) & $v$ \\
\code{linear} & \code{offset} (s/a, required), \code{slope\_x}, \code{slope\_y}, \code{slope\_z} (s/a, default 0) & $v_0+s_x x+s_y y+s_z z$ \\
\code{box} & \code{min\_corner}, \code{max\_corner} (2--3 numbers, min $\leq$ max), \code{inside\_values}, \code{outside\_values} (s/a) & inside value within the box, outside value elsewhere \\
\code{circle} & \code{center} (2--3 numbers), \code{radius} ($>0$), \code{dr} ($\geq 0$, required), \code{inside\_values}, \code{outside\_values} & disc (2-D) / sphere (3-D); $dr=0$ hard step, $dr>0$ blends inside $\to$ outside over [\code{radius}, \code{radius}$+dr$] with a $C^1$ cubic \\
\code{ellipse} & \code{center}, \code{semi\_axes} ($>0$ each), \code{inside\_values}, \code{outside\_values}, \code{smoothness} ($\geq 0$, default 0) & ellipse (ellipsoid) indicator; \code{smoothness} $=0$ hard step, $>0$ a $\tanh$ blend across the level set \\
\code{gaussian} & \code{center}, \code{sigma} ($>0$ each, per axis), \code{amplitude} (s/a), \code{offset} (s/a, default 0) & $v_0 + a\exp\bigl[-\sum_k (x_k-\mu_k)^2/(2\sigma_k^2)\bigr]$ \\
\code{cosine} & \code{wavenumber} (per axis), \code{phase} (per axis, default 0), \code{amplitude} (s/a), \code{offset} (s/a) & $v_0 + a\prod_k \cos(k_k x_k + \varphi_k)$ \\
\code{random} & \code{min}, \code{max} (min $<$ max), \code{seed} (required, non-negative integer) & i.i.d.\ uniform values in $[\min,\max]$ from a seeded \code{mt19937\_64} -- runs are bit-reproducible \\
\bottomrule
\end{tabularx}
\end{table}

A typical spinodal-decomposition start is a seeded random
perturbation about the mean composition:

\begin{lstlisting}[language=json]
"ICSystem": {
  "noise": { "type": "random", "dof": "c",
             "min": 0.49, "max": 0.51, "seed": 12345 }
}
\end{lstlisting}

%==================================================================
\section{Job Control, Time Stepping and Solvers}
%==================================================================

\subsection{\code{JobSystem}}
\label{sec:job}

Optional; when omitted the job is static with serial assembly. When
present, \code{type} is required.

\begin{table}[htbp]
\centering
\caption{\code{JobSystem} keys.}
\footnotesize
\begin{tabularx}{\textwidth}{lllX}
\toprule
\textbf{Key} & \textbf{Type} & \textbf{Default} & \textbf{Meaning} \\
\midrule
\code{type} & string & \code{"static"} & \code{"static"} (one solve) or \code{"transient"} (time loop; requires \code{TimeStepping}) \\
\code{assemble} & string & \code{"serial"} & \code{"serial"}, \code{"openmp"} (needs \code{PARALLEL\_ASSEMBLE}; thread count from \code{OMP\_NUM\_THREADS}) or \code{"cuda"} (needs \code{CUDA\_ASSEMBLE}) \\
\bottomrule
\end{tabularx}
\end{table}

\subsection{Run drivers}

The problem driver picks one of four run modes:

\begin{itemize}[nosep]
  \item \textbf{static} + implicit kernel: a single Newton solve,
        then one result save;
  \item \textbf{transient} + implicit kernel: the backward-Euler
        time loop of Section~\ref{sec:ts}, one Newton solve per
        step;
  \item \textbf{transient} + \code{explicit\_pddo\_frac}: the
        matrix-free central-difference loop
        (Section~\ref{sec:explicit});
  \item quasi-static matrix-free problems use an ADR
        (adaptive dynamic relaxation) driver -- not reachable with
        the published kernel set, listed for completeness.
\end{itemize}

\subsection{\code{TimeStepping}}
\label{sec:ts}

Required for transient jobs, forbidden for static ones.

\begin{table}[htbp]
\centering
\caption{\code{TimeStepping} keys.}
\footnotesize
\begin{tabularx}{\textwidth}{lllX}
\toprule
\textbf{Key} & \textbf{Type} & \textbf{Default} & \textbf{Meaning} \\
\midrule
\code{dt} & number & required & initial (and, non-adaptive, constant) step; $>0$ \\
\code{total\_time} & number & required & end time; $>0$ \\
\code{verbose} & bool & \code{true} & per-step progress lines \\
\code{adaptive} & bool & \code{false} & enable adaptive step control \\
\code{optimal\_iters} & int & 5 & Newton-iteration target of the controller; $\geq 1$ \\
\code{growth\_factor} & number & 1.25 & step multiplier after an easy step; $\geq 1$ \\
\code{cutback\_factor} & number & 0.5 & step multiplier after a failed step; $\in(0,1)$ \\
\code{max\_cutbacks} & int & 4 & retries per step before giving up; $\geq 0$ \\
\code{min\_dt} & number & 0 & lower bound on the step \\
\code{max\_dt} & number & 0 & upper bound; 0 means ``capped at the initial \code{dt}'' \\
\bottomrule
\end{tabularx}
\end{table}

The adaptive controller works per step: if Newton \emph{fails}, the
step is retried with $\Delta t \leftarrow \Delta t\cdot
\text{cutback\_factor}$, up to \code{max\_cutbacks} times, then the
run aborts. If Newton \emph{converges}, the next step grows by
\code{growth\_factor} when the iteration count was at most
\code{optimal\_iters} and shrinks by \code{cutback\_factor}
otherwise, always clamped to $[\text{min\_dt},\text{max\_dt}]$.
The last step is shortened to land exactly on
\code{total\_time}.

\subsection{\code{NonlinearSolver}}
\label{sec:newton}

Newton's method without line search: assemble
$\mathbf K\,\Delta\mathbf U=-\mathbf R$, apply boundary conditions,
solve, update $\mathbf U\mathrel{+}=\Delta\mathbf U$. Convergence
when $\lVert\mathbf R\rVert<\text{abs\_tol}$ or
$\lVert\mathbf R\rVert\le\text{rel\_tol}\cdot
\lVert\mathbf R_0\rVert$.

\begin{table}[htbp]
\centering
\caption{\code{NonlinearSolver} keys.}
\footnotesize
\begin{tabularx}{\textwidth}{lllX}
\toprule
\textbf{Key} & \textbf{Type} & \textbf{Default} & \textbf{Meaning} \\
\midrule
\code{max\_iters} & int & 50 & maximum Newton iterations per solve; $\geq 1$ \\
\code{abs\_tol} & number & $10^{-9}$ & absolute residual-norm tolerance; $>0$ \\
\code{rel\_tol} & number & $10^{-12}$ & relative residual-norm tolerance; $>0$ \\
\code{verbose} & bool & \code{true} & per-iteration residual lines \\
\bottomrule
\end{tabularx}
\end{table}

\subsection{\code{LinearSolver}}
\label{sec:linsolver}

\begin{table}[htbp]
\centering
\caption{Linear solver backends. Only the listed \code{params} keys
are public; anything else is rejected.}
\footnotesize
\begin{tabularx}{\textwidth}{llX}
\toprule
\textbf{\code{type}} & \textbf{Availability} & \textbf{Description and \code{params}} \\
\midrule
\code{default} & always & built-in direct profile (skyline) LDU factorisation without pivoting, with reverse Cuthill--McKee reordering to compress the profile. No parameters (\code{params} must be empty/omitted). \\
\code{pardiso} & \code{USE\_ONEAPI} builds & Intel MKL PARDISO sparse direct solver. Params: \code{msglvl} (int $\geq 0$, default 0) -- PARDISO statistics verbosity. \\
\code{cudss} & \code{USE\_CUDSS} builds & NVIDIA cuDSS sparse direct solver on the GPU. No public parameters. \\
\code{amgcl} & \code{USE\_AMGCL} builds & AMGCL BiCGSTAB with a preconditioner ladder (see below). Params: \code{tol} ($>0$, default $10^{-10}$), \code{maxiter} (int $\geq 1$, default 200), \code{verbose} (bool, default \code{false}), \code{iluk\_level} (int 1--6, default 2), \code{solver\_threads} (int $\geq 0$, default 0 = automatic). \\
\bottomrule
\end{tabularx}
\end{table}

The block itself is optional (the built-in direct solver is the
default). The AMGCL backend detects the block size (1--5 coupled
DoFs per node) automatically and climbs a preconditioner ladder --
AMG-with-ILU(0) relaxation first, plain ILU(0) next, ILU($k$) with
\code{iluk\_level} fill last -- restarting the Krylov solve when a
cheaper preconditioner stalls. Requesting \code{pardiso},
\code{cudss} or \code{amgcl} in a binary compiled without the
corresponding CMake option is a start-up error.

%==================================================================
\section{Output System}
\label{sec:output}
%==================================================================

\subsection{The \code{OutputSystem} block}

\begin{table}[htbp]
\centering
\caption{\code{OutputSystem} keys.}
\footnotesize
\begin{tabularx}{\textwidth}{lllX}
\toprule
\textbf{Key} & \textbf{Type} & \textbf{Default} & \textbf{Meaning} \\
\midrule
\code{interval} & int & 1 & save every $n$-th step ($\geq 1$); step 0 and the final step are always saved \\
\code{Fields} & array & \code{[]} & projection names to write in addition to the solution; duplicates are rejected; \code{"solution"} is accepted and ignored (the solution is always written); a name no kernel provides is a fatal error \\
\code{format} & string & \code{"vtu"} & \code{"vtu"}, \code{"exodus"} or \code{"both"} \\
\bottomrule
\end{tabularx}
\end{table}

Static jobs save once after the solve. The available projection
names per kernel are listed in Section~\ref{sec:elmt}; vector and
tensor projections are written with their natural number of
components.

\subsection{VTK output}

Two files (per saved step) plus their ParaView index:

\begin{itemize}[nosep]
  \item \code{<stem>-PD-results[-NNNNN].vtu}: the PD point cloud as
        vertex cells, with point data \code{ElmtID}, one scalar per
        DoF (canonical names \code{u}, \code{c}, \code{mu},
        \code{ux}, \ldots) and the requested projections;
  \item \code{<stem>-Element-results[-NNNNN].vtu}: the background
        mesh with the same data as cell data (one PD point per
        cell), convenient for contour plots;
  \item \code{<stem>-PD-results.pvd} /
        \code{<stem>-Element-results.pvd}: time-series indexes,
        rewritten at every save so a running simulation can be
        opened at any moment.
\end{itemize}

Transient file names carry a five-digit zero-padded step counter;
static results have no counter.

\subsection{ExodusII output}

With \code{"format": "exodus"} (or \code{"both"}) \PeriX{} writes
\code{<stem>.e} -- a self-contained ExodusII database in CDF-2
format produced by the built-in writer; \emph{no} NetCDF or SEACAS
libraries are needed at build or run time. The mesh is written once
with the proper element types (\code{BAR2}, \code{TRI3},
\code{QUAD4}, \code{TETRA4}, \code{HEX8}); every save appends a
time plane with strictly increasing times. Each DoF becomes an
element variable under its canonical name; multi-component
projections are split into \code{name\_1}, \code{name\_2}, \ldots{}
components.

%==================================================================
\section{Worked Examples}
\label{sec:examples}
%==================================================================

The \code{examples/} directory holds eleven decks covering every
kernel, every mesh type, every solver backend and both assembly
accelerators. Each deck is schema-validated by the test suite.
Run any of them with

\begin{lstlisting}[language=bash]
cd examples
../build/perix -i poisson.json
\end{lstlisting}

\subsection{\code{poisson.json} -- steady conduction}

A $1000\times1000$ unit square ($10^6$ PD points), horizon factor
3.0, order 2, conductivity $\sigma=1$, no source. Dirichlet walls
$u=0$ (left) and $u=1$ (right) through the mirror form on the ghost
groups; top and bottom carry no condition. OpenMP assembly, PARDISO
solver, one static solve. The complete deck:

\begin{lstlisting}[language=json]
{
  "Mesh":   { "type": "perix", "dim": 2, "xmax": 1.0, "ymax": 1.0,
              "nx": 1000, "ny": 1000, "meshtype": "quad4",
              "savemesh": false },
  "PDMesh": { "HorizonRadiusFactor": 3.0, "Order": 2 },
  "DOFs":   [ "u" ],
  "ElmtSystem": {
    "poisson": { "type": "poisson",
                 "params": { "sigma": 1.0, "f": 0.0 } }
  },
  "BCSystem": {
    "left":  { "type": "dirichlet", "phygroup": "leftnodes_ghost",
               "dofs": ["u"], "value": 0.0 },
    "right": { "type": "dirichlet", "phygroup": "rightnodes_ghost",
               "dofs": ["u"], "value": 1.0 }
  },
  "JobSystem":    { "type": "static", "assemble": "openmp" },
  "LinearSolver": { "type": "pardiso" }
}
\end{lstlisting}

The solution is the linear ramp $u=x$; a useful first check of any
build.

\subsection{\code{poisson\_amgcl.json} and \code{poisson\_gpu.json}}

The same physics on a $50\times50$ grid, demonstrating the optional
backends: \code{poisson\_amgcl.json} solves with
\code{"type": "amgcl"} (\code{tol} $10^{-10}$, \code{maxiter} 200)
and serial assembly; \code{poisson\_gpu.json} uses
\code{"assemble": "cuda"} with the \code{cudss} direct solver.
They require binaries configured with \code{USE\_AMGCL} and
\code{CUDA\_ASSEMBLE}+\code{USE\_CUDSS} respectively.

\subsection{\code{diffusion.json} -- transient diffusion}

A $500\times500$ unit square, $D=0.01$, $c$ initially zero,
Dirichlet $c=1$ at the bottom wall and $c=0$ at the top, zero-flux
Neumann on the sides. Backward Euler with fixed
$\Delta t=0.5$ to $t=50$; results every 10 steps. The concentration
front diffuses upward towards the steady linear profile.

\subsection{\code{spinodal.json} -- Cahn--Hilliard phase separation}

A $250\times250$ unit square with $\chi=3$ (deep quench),
$\kappa=10^{-3}$, $M_0=1$. The initial state is the seeded random
perturbation $c\in[0.49,0.51]$ (\code{seed} 12345, bit-reproducible)
with $\mu=0$; all four walls are homogeneous Neumann. The adaptive
controller (initial $\Delta t=10^{-3}$, \code{optimal\_iters} 10,
growth 1.2, cutback 0.5, \code{max\_cutbacks} 6, $\Delta t$
clamped to $[10^{-4},1]$) rides the stiffness of early spinodal
decomposition and accelerates through late-stage coarsening to
$t=100$. Fields \code{fPrime} and \code{fDoublePrime} are written
every 20 steps; total mass $\sum_i V_i c_i$ is conserved to solver
tolerance.

\subsection{\code{kalthoff\_winkler.json} -- explicit dynamic fracture}

The classic Kalthoff--Winkler impact test: a pre-notched steel
plate (Gmsh mesh \code{impact2D.msh}) struck between the notches.
Maraging-steel parameters $E=190\,$GPa, $\nu=0.25$,
$\rho=7800\,$kg/m$^3$, $G_0=69\,000\,$J/m$^2$, plane stress,
tension-only failure. The two notches are explicit \code{Cracks}
segments at $x=0.075$ and $x=0.125$ (from $y=0.05$ to $y=0.1$) with
\code{force\_only} treatment. The impactor region
(\code{impactnodes\_ghost}) is driven by a direct Dirichlet
condition with \code{velocity} $[0,-22]$\,m/s; the remaining
boundary (\code{freenodes\_ghost}) is traction-free via
homogeneous Neumann. Central-difference integration with
$\Delta t=10^{-8}\,$s to $t=10^{-4}\,$s; \code{damage},
\code{vonMises} and \code{stress} saved every 1000 steps. The
simulation reproduces the experimentally observed crack branches at
roughly $70^\circ$ to the notch line.

\code{kalthoff\_winkler\_gpu.json} is the same deck with
\code{"assemble": "cuda"} -- the entire bond loop runs on the GPU.

\subsection{\code{kalthoff\_winkler\_3D.json} -- 3-D explicit fracture}

The 3-D extension on \code{impact3D.msh} with DoFs
\code{ux},\code{uy},\code{uz} and \code{Gc} spelled directly
(no \code{state} in 3-D). The notches are \code{center\_crack}
presets with \code{normal} $[1,0,0]$, \code{length} 0.05 and
\code{width} 0.009 -- planar cracks through the plate thickness.
The out-of-plane displacement \code{uz} is pinned on
\code{alldomain} to keep the thin plate in a plane-strain-like
state, and the impactor drives \code{velocity} $[0,-22,0]$.
$\Delta t=5\times10^{-8}\,$s, output every 100 steps.

\subsection{\code{tensile\_plate.json} -- implicit quasi-static fracture}

A $1\times0.4$ plate (300$\times$120 cells) with an edge crack from
$(0,0.2)$ to $(0.5,0.2)$, loaded by \code{pdtraction} conditions:
$\pm0.0153884$ vertical traction on the top/bottom walls and
traction-free left/right walls (all with $E=79.1403$,
$\nu=0.333$, plane stress -- a normalised material set,
$\rho=268.198$, $G_c=1.48388\times10^{-6}$). Note that the traction
conditions address the \emph{bulk} wall groups
(\code{topnodes}, \ldots); the condition automatically maps each
bulk node to its mirror ghost and writes the traction rows there.
The implicit \code{pddo\_dynamic\_frac} kernel advances with the
adaptive controller ($\Delta t$ from $10^{-3}$ in $[10^{-4},10^{-2}]$,
growth 1.2, cutback 0.5) to $t=5$; the crack propagates across the
ligament and the plate tears in two.

\subsection{\code{silicon.json} -- lithiation of a silicon disc}

Stress-assisted Cahn--Hilliard on the \code{circle} mesh
(radius 1, 400 layers, \code{VariableHorizon} on): a normalised
model of Li insertion into an amorphous-Si particle. Parameters
$E=87.9336$, $\nu=0.29$, $D=10^{-4}$, $\Omega=3$,
$c_{\mathrm{ref}}=0$, $\chi=2.6$,
$\kappa=2.19834\times10^{-6}$, $\rho=2.41817\times10^{-22}$
(inertia negligible but nonzero), plane stress. Boundary
conditions on the rim (\code{outernodes\_ghost}): a constant
species influx $j=0.013889$ (\code{speciesflux}), zero-gradient
\code{mu}, and a swelling-aware traction-free surface
(\code{sdtraction} with the same $E,\nu,\Omega$). A small box
around the centre pins \code{ux},\code{uy} to remove rigid-body
modes. Adaptive stepping to $t=30$; \code{stress} and
\code{sigmaH} written every 50 steps. The concentration front
sweeps inward while the core is pushed into hydrostatic
compression -- the classical two-phase lithiation picture.

\subsection{\code{silicon\_frac.json} -- lithiation with a pre-crack}

The fracture variant: \code{frac\_stress\_cahnhilliard} with
$G_c=0.010662$, \code{tension\_only} true, faster transport
($D=0.03$) and a radial pre-crack from $(-2,0)$ to $(-0.5,0)$ --
the segment reaches beyond the disc so the crack terminates
exactly at the rim. With \code{"damage\_on": false} the bond
network stays frozen: the deck isolates the effect of the
pre-crack (stress concentration, locally re-routed swelling
stresses) from crack \emph{growth}, which can be enabled by
flipping the switch to \code{true}.

%==================================================================
\section{Testing}
%==================================================================

The test tree registers 29 CTest cases under four (overlapping)
labels:

\begin{itemize}[nosep]
  \item \textbf{unit} (17 tests): one test executable per
        subsystem -- every element kernel, the BC/IC systems, the
        circle-mesh anchor, the Exodus writer, initial velocities,
        the job system, \code{MeshModify}, the output system, the
        input schema, and the AMGCL solver components;
  \item \textbf{examples} (13): the input-schema suite plus a
        \code{ReadOnly\_<deck>} validation of each of the eleven
        bundled decks (\code{perix -i deck --read-only}), so every
        shipped input is guaranteed to parse and validate;
  \item \textbf{integration} (1): \code{Run\_poisson\_amgcl}, a
        full end-to-end solve of the AMGCL Poisson deck with
        checked results;
  \item \textbf{amgcl} (2): the AMGCL unit suite and the
        end-to-end run (present only in \code{USE\_AMGCL} builds).
\end{itemize}

\begin{lstlisting}[language=bash]
cd build && ctest -j 16 --output-on-failure
\end{lstlisting}

%==================================================================
\section{License and Citation}
%==================================================================

\PeriX{} is distributed under the GNU General Public License v3.0.
The PDDO method is due to Madenci and co-workers; the standard
references are:

\begin{itemize}[nosep]
  \item E.~Madenci, A.~Barut, M.~Futch,
        \emph{Peridynamic differential operator and its
        applications}, Computer Methods in Applied Mechanics and
        Engineering 304 (2016) 408--451.
  \item E.~Madenci, A.~Barut, M.~Dorduncu,
        \emph{Peridynamic Differential Operator for Numerical
        Analysis}, Springer, 2019.
\end{itemize}

Please cite the repository
(\url{https://github.com/MatMechLab/PeriX}) when using \PeriX{} in
published work.

\end{document}
